\subsection{The actual three-step inverse-$q$ fourth moment}\label{proof:Q5-C056}
We retain Q5-C056, Section 3, with its necessary elementary premises.
On original positive-mass support $q(x)\ge\mu(x)>0$.
Use $r=Kg$, $s=Kr$, $d=\E g^2$, $\delta_2=\E r^2$,
$\delta_3=\E s^2$.
The exact original skew identities are
$\E g=0$, $\E r=-d$, $\E s=0$,
$\E gr=\E rs=0$, $\E gs=-\delta_2$.
The actual paths are $\bar f_1=q$, $\bar f_2=q+r$,
$\bar f_3=q+r-s-d$.
We need only the elementary weaker spectral bound $\|K\|^2\le1/2$:
the original Hilbert--Schmidt square is at most one; for a maximizing
real unit vector $v$, $u=Kv/\|K\|$ is orthogonal to $v$ and
$\|Ku\|\ge\|K\|$. Extending $u,v$ to an orthonormal basis gives
$\|K\|_{\rm HS}^2\ge2\|K\|^2$.
Thus $\delta_2\le d/2$, $\delta_3\le\delta_2/2$.
The original three-cycle inequality gives $d\le1/3$.

Put $v=r-s$, $u_0=v-d=\bar f_3-q$, $W=2v-g$.
Their exact moments are
\[
\begin{gathered}
\E v=-d,\quad \E gv=\delta_2,\quad \E v^2=\delta_2+\delta_3,\quad \E u_0=-2d,\quad \E gu_0=\delta_2,\\
\E u_0^2=\delta_2+\delta_3+3d^2,
\end{gathered}
\]
\[
\E W=-2d,\qquad\E W^2=d+4\delta_3.
\]
The latter equality keeps the cancellation
$4(\delta_2+\delta_3)-4\delta_2+d$.
The actual generating identity is $W=K(-1+2g-2r)$,
with input square $1+4\delta_2$. Therefore the original row bound
gives $|W|\le\sqrt{1+4\delta_2}\le\sqrt{1+2d}\le1+d=:t$.
This includes small-mass atoms where $|g|$ is large.
The correlated cubic combination is
\[
\mathcal C=-\E g^3+4\E(g^2v)-6\E(gv^2)+4\E v^3
=\tfrac12\{\E W^3+\E(g^2W)\}.
\]
Since $(W+2d)^2(W-t)\le0$, exact expansion and
$\E W=-2d$ give
$\E W^3\le(t-4d)\E W^2-4td^2+8d^3$.
Also $\E(g^2W)\le td$, hence
\begin{equation}\tag{L56.3.8}\label{eq:L56:cubic}
\mathcal C\le td-2(1+t)d^2+4d^3+2(t-4d)\delta_3.
\end{equation}
The cubic interpolation applies for every $W$.
For the only inverse-weight quartic term, use the actual identity
$u_0=T(\bar f_2-1)$.
For $w=\bar f_2-1$, $\E w^2=d+\delta_2$;
row Cauchy--Schwarz and the original kernel bound two give
$u_0^2\le qT(w^2)$, $T(w^2)\le2(d+\delta_2)$.
Thus
\begin{equation}\tag{L56.3.9}\label{eq:L56:inverse}
\E(u_0^4/q)\le2(d+\delta_2)(\delta_2+\delta_3+3d^2).
\end{equation}
This controls arbitrarily small original $q$ without a uniform bound
on $1/q$.
The exact fourth expansion is
$R_4:=\E(\bar f_3^4/q)
=\E q^3+4\E(q^2u_0)+6\E(qu_0^2)+4\E u_0^3+\E(u_0^4/q)$,
with individual terms
\[
\begin{aligned}
\E q^3&=1+3d-\E g^3,\\
\E(q^2u_0)&=-2d-2\delta_2+\E(g^2v)-d^2,\\
\E(qu_0^2)&=\delta_2+\delta_3+3d^2-\E(gv^2)+2d\delta_2,\\
\E u_0^3&=\E v^3-3d(\delta_2+\delta_3)-4d^3.
\end{aligned}
\]
The $d\delta_2$ terms cancel exactly, leaving
$R_4=1-5d-2\delta_2+6\delta_3+14d^2-16d^3
-12d\delta_3+\mathcal C+\E(u_0^4/q)$.
Combining \eqref{eq:L56:cubic}--\eqref{eq:L56:inverse} gives
\[
R_4-1\le-(5-t)d-2\delta_2+(6+2t-20d)\delta_3
+(12-2t)d^2-12d^3
+2(d+\delta_2)(\delta_2+\delta_3+3d^2).
\]
For $d>0$ put $x=\delta_2/d\in[0,1/2]$ and substitute $t=1+d$.
The separate $\delta_3$ coefficient is $8-18d\ge2$ and its
coefficient in the last product is also nonnegative, so
$\delta_3\le\delta_2/2$ can be inserted legally everywhere.
The normalized upper bound becomes
\[
-4+11d+(2-6d+6d^2)x+3dx^2-8d^2.
\]
Its $x$ derivative is positive since
$2-6d+6d^2=6(d-1/2)^2+1/2>0$.
At $x=1/2$ it is $-3+(35/4)d-5d^2$.
This increases throughout $[0,1/3]$, derivative at least $65/12$,
and at $1/3$ is exactly $-23/36$.
We obtain the complete original-measure interface
\begin{equation}\tag{L56.3.15}\label{eq:L56:claim}
\E[\bar f_3^4/q]\le1-23d/36.
\end{equation}
At $d=0$, $g=r=s=0$, all these actual path messages and weights are
one, while $K$ may be nonzero. The proof uses no lower spectral
moment inequality or stronger spectral constant.
